\clearpage
\VerdierPhysicalPageAnchor{0228}%
\begingroup\small\setlength{\parskip}{0pt}\setlength{\emergencystretch}{1.5em}
\setlength{\abovedisplayskip}{4pt plus 2pt minus 2pt}%
\setlength{\belowdisplayskip}{4pt plus 2pt minus 2pt}%

\VerdierSemanticAnchor{ch4d-4-1-4}
\medskip\noindent\textbf{4.1.4.}\quad
Il résulte de (4.1.3) que pour tout objet $Y$ de
$\operatorname{comp}(\mathcal A)$,

\VerdierSemanticAnchor{ch4d-4-1-4-spec-k}
\[
\operatorname{Spec}_1(Y)=\bigl(Y(p,q),\delta(p,q,r)\bigr),
\]
défini par (4.1.2.3), (4.1.2.4) et (4.1.2.6), est un objet spectral à
valeurs dans $K(\mathcal A)$. L'objet spectral $\operatorname{Spec}_1(Y)$
dépend fonctoriellement de l'objet $Y$ de
$\operatorname{comp}(\mathcal A)$. On a donc déterminé un foncteur :

\VerdierSemanticAnchor{ch4d-4-1-4-functor}
\[
Y\longmapsto\operatorname{Spec}_1(Y)
\]
de la catégorie des complexes de $\mathcal A$ dans la catégorie des objets
spectraux à valeurs dans $K(\mathcal A)$. Ce foncteur est appelé le
\emph{foncteur premier objet spectral canonique à valeurs dans
$K(\mathcal A)$}. Pour tout objet $Y$ de
$\operatorname{comp}(\mathcal A)$, l'objet spectral
$\operatorname{Spec}_1(Y)$ est appelé le \emph{premier objet spectral à
valeurs dans $K(\mathcal A)$ associé au complexe $Y$}. Il convient de noter
que lorsque $Y$ est un objet de $\operatorname{comp}^{*}(\mathcal A)$, où
$*$ vaut $+$, $-$, $b$ ou est ``vide'', le premier objet spectral à valeurs
dans $K(\mathcal A)$ associé au complexe $Y$ est, en fait, à valeurs dans
$K^{*}(\mathcal A)$.

\VerdierSemanticAnchor{ch4d-4-1-5}
\medskip\noindent\textbf{4.1.5.}\quad
Soient $\mathcal A$ une catégorie abélienne et $Y$ un complexe de
$\mathcal A$. L'image par le foncteur canonique
$K(\mathcal A)\longrightarrow D(\mathcal A)$ du premier objet spectral
associé au complexe $Y$ est un objet spectral à valeurs dans
$D(\mathcal A)$ qui est appelé le \emph{premier objet spectral à valeurs
dans $D(\mathcal A)$ associé au complexe $Y$}. Le plus souvent, nous dirons
simplement \emph{premier objet spectral associé au complexe $Y$} et
noterons (abusivement) :

\VerdierSemanticAnchor{ch4d-4-1-5-spec-d}
\[
\operatorname{Spec}_1(Y)=\bigl(Y(p,q),\delta(p,q,r)\bigr)
\]
le premier objet spectral à valeurs dans $D(\mathcal A)$ associé au
complexe $Y$. Le foncteur premier objet spectral canonique composé avec le
foncteur canonique $K(\mathcal A)\longrightarrow D(\mathcal A)$ (étendu aux
objets spectraux) est appelé le \emph{foncteur premier objet spectral
canonique à valeurs dans $D(\mathcal A)$}. Soulignons qu'il s'agit d'un
foncteur de la catégorie des complexes de $\mathcal A$ dans la catégorie
des objets spectraux de $D(\mathcal A)$ et que \emph{ce foncteur ne se
factorise pas, en général, par la catégorie des complexes à homotopie près}.

\VerdierSemanticAnchor{ch4d-4-1-6}
\medskip\noindent\textbf{4.1.6.}\quad
Les complexes $n$-uples d'objets de $\mathcal A$ permettent de construire
des complexes simples munis de filtrations et par suite des objets spectraux
à valeurs dans $K(\mathcal A)$. Soient $X$ un complexe $n$-uple (chap. I,
2.1.2) et $m$ un entier, $1\leq m\leq n$. Pour tout
$\sigma\in\mathbb Z[n]$ (notation de (chap. I, 2.1.1)), désignons par
$\sigma_m$ sa $m$-ième coordonnée.

\VerdierSemanticAnchor{ch4d-4-1-6-family-setup}
Nous allons, tout d'abord, définir une famille de complexes et de morphismes
de complexes. Soit, pour tout couple $(p,q)$ de
$\mathcal{F}\ell_1(\widetilde{\mathbb Z})$, ${}^{m}X(p,q)$ le
\endgroup
